\section{Relation to positive realization and stochastic automata}\label{sec:literature45}
For fixed seed and query, the centered response is a rational word series of signed linear dimension at most $D$. Positivity, normalization, and common continuation are additional requirements, not consequences of that dimension. The following comparisons locate those requirements before drawing any conclusion from the arithmetic examples.

\subsection{Signed Hankel rank and positive order}
Balle--Panangaden--Precup \cite[Theorem 2]{BPP} recall the classical equivalence between finite Hankel rank and minimum weighted-automaton dimension. It applies to signed rational series. In our notation it explains the first rank bound preceding \eqref{eq:rank-chain45}, not the last inequality in the opposite direction. Their SVA truncation result \cite[Theorem 11]{BPP} concerns $\ell^2$ error for absolutely summable rational series under a stability hypothesis. Our undamped orthogonal word series need not satisfy that hypothesis, and the approximating rows here must remain stochastic and all terminal columns bounded. A finite-length $\ell^2$ error would bound a finite-length supremum error, but these positivity and stability conditions would still have to be established. We do not infer a stochastic compiler from signed minimization.

\subsection{Positive cones and the ambient space}
Benvenuti--Farina \cite[Theorems 2--3]{BF} formulate positive realization through invariant polyhedral cones, including the shift representation. Their Section VII also explains why restricting a cone to a prescribed minimal observable realization can fail to recover the unrestricted minimum positive order. Theorem~\ref{thm:flag45} is a finite, multi-command, normalized version of that mechanism. Its entire future-response cube is essential; replacing it by the $D$-dimensional target space would silently restrict the competing hidden states. The enclosure of Theorem~\ref{thm:enclosure} is a useful special cone section, not all such flags. Benvenuti's later survey \cite{Ben22} provides further background on minimum positive order.

Stationary realization asks for one unchanged set of transition and output data at every length. The present quantifiers are instead ``for every $N$ there exists a layered machine'', with the final decoder depending on $N$. The one-letter irrational example is a concrete distinction, rather than an assertion that all classical stationary lower bounds transfer. Likewise Proposition~\ref{prop:rank-three45} shows that even every separate positive factorization can have size three while the compatible full-run width diverges.

\subsection{Stochastic realization and factorization consistency}
Heller's finite stochastic-realization framework \cite{Heller} is part of the classical lineage of shift-invariant positive cones. Ohta \cite[Theorems 6--9 and Remark 10]{Ohta} analyzes reachable and null spaces, effective dimension, tensor decomposition and generalized Hankel matrices for stationary hidden Markov models. These distinguish observable dimension from a particular hidden presentation. In our controlled model, each command is externally specified and each row is stochastic separately; an emission-indexed HMM kernel instead has a different normalization across its emissions. Neither a stationary effective-space reduction nor a tensor-rank bound alone supplies \eqref{eq:flag45} for every command and epoch.

Finesso--Grassi--Spreij \cite[Section IV]{FGS} explicitly separate pseudo-Hankel nonnegative factorization from the subsequent parametrization step needed to obtain HMM transition data. Their objective uses informational divergence of finite distributions. This is directly relevant to the warning after \eqref{eq:rank-chain45}: independent factors are not yet a common stochastic dynamics. Our exact flag formula uses worst-word binary total variation, equivalent to one half the supremum error of means. It is a finite characterization, not a convex solution of the generally nonconvex factorization problem.

\subsection{What the boundedness criterion adds}
The coset obstruction in group random walks is classical; for the finite-group aperiodicity statement, see the introduction of McCarthy \cite{McCarthy}. The subgroup in \eqref{eq:sync-intro} is the smallest closed normal subgroup in whose coset all commands lie. We do not claim that this group-theoretic object, invariant cones, or finite-group recognition was first introduced here. Our conclusion is the equivalence of its finiteness on the active space with uniformly bounded horizon-specific \emph{numerical} positive width. Its converse combines equal-length orbit stabilization, compact uniformization on the moving subspace, and the inherited arbitrary-hidden-state packet inequality. It applies without a gap, with noncommuting commands, and across disconnected actions.

Finite matrix-group recognition is an established algorithmic subject. Detinko--Flannery--O'Brien \cite[Sections 3.1--3.2 and 4]{DFO} give much stronger computational methods than our coarse enumeration. Theorem~\ref{thm:decision45} uses a proved order bound to close the normal-conjugation enumeration and turns the new memory criterion into a finite algebraic decision. It is neither a new general finiteness algorithm nor an infinite spectral-gap oracle.

The arithmetic results below retain their original distinctions: fixed alphabet size, fixed calibration, and logarithmic exponents rather than constant-factor laws where only ordinary Diophantine type is assumed. Their standard individual versus uniform exponent convention is made explicit using German \cite[Definitions 1--2]{German}: the former asks for infinitely many approximations, whereas the latter asks for one at every sufficiently large scale. The row-vector dual exponent in \eqref{eq:type-intro} is the individual exponent in that convention; changing from the supremum norm to $\ell^1$ does not change its value. The direct DGLPS comparison is retained in the final arithmetic discussion.
